\newpage
% Оформляем как обязательное приложение (например, К)
\begin{flushright}
    {\bfseries ПРИЛОЖЕНИЕ К} \\
    (обязательное)
\end{flushright}

\begin{center}
    {\bfseries ТЕХНИКО-ЭКОНОМИЧЕСКИЕ ПОКАЗАТЕЛИ ПРОЕКТА}
\end{center}

\addcontentsline{toc}{section}{Приложение К. Технико-экономические показатели}

\vspace{0.5cm}

\noindent В таблице К.1 приведены основные результаты расчета экономической эффективности разработанной системы мониторинга в сравнении с базовым вариантом (аналогом).

\vspace{0.3cm}

\begin{table}[h!]
    \centering
    \caption*{Таблица К.1 --- Сводные технико-экономические показатели}
    \begin{tabular}{|p{7cm}|c|c|c|}
        \hline
        \centering \textbf{Наименование показателя} & \textbf{Ед. изм.} & \textbf{Аналог} & \textbf{Проект} \\ \hline
        1. Затраты на разработку (КВ) & руб. & --- & 5697,20 \\ \hline
        2. Годовые эксплуатационные затраты & руб. & 1240,00 & 450,00 \\ \hline
        3. Потребляемая мощность & Вт & 15,0 & 2,5 \\ \hline
        4. Средняя наработка на отказ (MTBF) & ч. & 10000 & 25000 \\ \hline
        5. Годовая экономия (эффект) & руб. & --- & 3200,00 \\ \hline
        6. Рентабельность затрат & \% & --- & 56,1 \\ \hline
        7. Срок окупаемости & лет & --- & 1,78 \\ \hline
    \end{tabular}
\end{table}

\vspace{1cm}

\noindent \textbf{Вывод:} Анализ данных таблицы К.1 показывает, что разработанная система превосходит существующий аналог по показателям энергоэффективности и надежности. Срок окупаемости проекта составляет менее двух лет, что подтверждает высокую экономическую целесообразность внедрения разработки на предприятии.

\vfill

\begin{flushleft}
    \begin{tabular}{p{6cm} p{4cm} l}
        Разработал (студент) & \underline{\hspace{3cm}} & П. П. Петров \\
        Консультант по экономике & \underline{\hspace{3cm}} & Е. Е. Экономистов
    \end{tabular}
\end{flushleft}

\newpage